\documentclass[12pt,a4paper]{article}

% -------------------------------------------------
% OSNOVNE NASTAVITVE
% -------------------------------------------------
\usepackage[utf8]{inputenc}
\usepackage[T1]{fontenc}
\usepackage[slovene]{babel}
\usepackage{lmodern}
\usepackage[a4paper,margin=2.5cm]{geometry}
\usepackage{setspace}
\onehalfspacing

\usepackage{amsmath,amssymb}
\usepackage{booktabs}
\usepackage{array}
\usepackage{tabularx}
\usepackage{graphicx}
\usepackage{float}
\usepackage{caption}
\usepackage{url}
\usepackage[hidelinks]{hyperref}
\usepackage{csquotes}

% Grafi neposredno iz CSV
\usepackage{pgfplots}
\pgfplotsset{compat=1.18}

% Literatura
\usepackage[
    backend=biber,
    style=authoryear,
    sorting=nyt,
    uniquename=false,
    uniquelist=false
]{biblatex}
\addbibresource{viri.bib}

% -------------------------------------------------
% PODATKI O NALOGI
% -------------------------------------------------
\title{\textbf{Umetna inteligenca kot vir R\&R in dolgoročne gospodarske rasti:}\\
ZDA v obdobju 2010--2025}
\author{
Andraž Ajdovec \and Tin Cvek \and Jan Hrovat \and Miha Ručigaj
}
\date{oktober 2026}

\begin{document}

\maketitle
\thispagestyle{empty}

\vfill
\begin{center}
Predmet: Makroekonomija\\
Seminarska naloga\\
Univerza v Ljubljani, Fakulteta za matematiko in fiziko\\
Mentor: \textit{Jan Grobovšek}
\end{center}
\vfill

\newpage
\tableofcontents
\newpage



% =================================================
\section{Uvod}\label{sec:uvod}
% =================================================

\subsection{Širši okvir in pomen problematike}
Gospodarska rast je razmeroma nov pojav, ki se je izrazito pospešil okoli leta 1800 v Zahodni Evropi in njenih izpeljankah, med katere spadajo tudi ZDA \parencite[pros.~3]{grobovsek2026}. Dinamiko vodilnih industrijskih držav, kot sta Združeno kraljestvo in ZDA, lahko v zadnjih dveh stoletjih v logaritmih približamo z ravno črto, kar pomeni približno konstantno povprečno stopnjo rasti \parencite[pros.~5]{grobovsek2026}. Že majhne razlike v tej stopnji so pomembne: ena odstotna točka višja rast v štiridesetih letih pomeni približno 49~\% višjo raven dohodka na prebivalca \parencite[str.~4]{polanec_rast}, zato so koristi višje dolgoročne rasti bistveno večje od koristi odprave kratkoročnih nihanj \parencite[str.~56]{garin2021}.

Osrednji model za razumevanje rasti je Solowov model. Njegova ključna implikacija je, da je dolgoročna stopnja rasti BDP na delavca enaka stopnji tehnološkega napredka in da je tehnologija edini vir dolgoročne rasti dohodka na delavca \parencite[pros.~59]{grobovsek2026}. Tehnologijo razumemo kot zalogo idej, ki so nerivalni proizvodni dejavnik; med primeri idej predavanje navaja tudi vrednosti parametrov v nevronski mreži \parencite[pros.~115--117]{grobovsek2026}. Ideje nastajajo z investicijami v raziskave in razvoj (R\&R) \parencite[pros.~124]{grobovsek2026}. Za ZDA ocene kažejo, da je družbena donosnost izdatkov za R\&R, ki znašajo okoli 2,7~\% BDP, zelo visoka (Jones in Summers, 2020, povzeto po \cite[pros.~131]{grobovsek2026}).

Umetna inteligenca (UI) je zato za teorijo rasti zanimiva iz dveh razlogov. Po eni strani nekateri ekonomisti, npr.\ Gordon (2016), ocenjujejo, da postaja vse težje ustvarjati nove prodorne ideje, saj je rast produktivnosti v ZDA kljub hitremu povečevanju števila raziskovalcev ostala približno nespremenjena (Bloom in sod., 2020, povzeto po \cite[pros.~136--137]{grobovsek2026}). Po drugi strani predavanje UI uvršča med makro inovacije, ki se šele zares začenjajo \parencite[pros.~139]{grobovsek2026}. Ocene njenega učinka se močno razlikujejo: \textcite[29]{acemoglu2024} na primer ocenjuje, da bo UI v naslednjih desetih letih letno rast TFP v ZDA zvišala le za okoli 0,06 odstotne točke.

\subsection{Predstavitev izbranega primera: ZDA v obdobju 2010--2025}
Za empirični del smo izbrali ZDA iz treh razlogov. Prvič, so gospodarstvo z največjim obsegom R\&R: skupni izdatki za R\&R so leta 2023 znašali 939,6 milijarde USD \parencite[tab.~DISC-1]{nsb2025}. Drugič, v ZDA so bile zasebne naložbe v UI leta 2024 skoraj dvanajstkrat večje kot na Kitajskem \parencite{aiindex2025}. Tretjič, ZDA so na tehnološki meji. Produktivnost gospodarstva lahko razčlenimo na svetovno mejno tehnologijo in učinkovitost, s katero jo gospodarstvo uporablja; prva je vir dolgoročne rasti, druga pa določa raven \parencite[pros.~119--120]{grobovsek2026}. Pri gospodarstvu na meji rast produktivnosti zato v veliki meri odraža rast mejne tehnologije.

Obdobje 2010--2025 zajame leta pred širšo uporabo UI in leta njenega hitrega razmaha; generativna UI se je množično razširila po izidu orodja ChatGPT konec leta 2022. Obdobje razdelimo na tri podobdobja (2010--2014, 2015--2019 in 2020--2025), zadnjega pa še na leti pandemije (2020--2022) in leta po izidu orodja ChatGPT (2023--2025).

\subsection{Raziskovalno vprašanje}\label{sec:rv}
\noindent\textbf{Raziskovalno vprašanje:} Ali je pospešek rasti produktivnosti v ZDA v obdobju 2010--2025 po Solowovem modelu posledica višjih investicij ali hitrejšega tehnološkega napredka in ali ga lahko pripišemo umetni inteligenci?

Vprašanje razdelimo na tri podvprašanja:
\begin{enumerate}
    \item Ali so se povečala vlaganja v R\&R in UI ter celotna stopnja investicij?
    \item Ali je razmerje med kapitalom in BDP ostalo konstantno, kot model napoveduje za uravnoteženo pot, in koliko rasti BDP na uro dela pojasnijo to razmerje, človeški kapital in produktivnost?
    \item Ali je pospešek rasti produktivnosti robusten in ali ga glede na časovni potek in velikost lahko pripišemo UI?
\end{enumerate}

Za odgovor uporabimo Solowov model z rastjo prebivalstva in tehnologije, ki ga v poglavju~\ref{sec:teorija} teoretično predstavimo, v poglavju~\ref{sec:uporaba} pa ga z računovodsko razčlenitvijo rasti iz predavanj uporabimo na podatkih. Vzročnosti ne dokazujemo, temveč preverjamo napovedi modela in alternativne razlage.

\subsection{Struktura naloge}
V poglavju~\ref{sec:podatki} predstavimo podatke o izdatkih za R\&R, investicijah, vlaganjih v UI, UI-patentih in produktivnosti v ZDA. V poglavju~\ref{sec:teorija} predstavimo predpostavke, enačbe in uravnoteženo rast Solowovega modela, njegovo računovodsko razčlenitev ter možne učinke UI s številčno simulacijo. V poglavju~\ref{sec:uporaba} model uporabimo na podobdobjih, v poglavju~\ref{sec:zakljucek} pa odgovorimo na raziskovalno vprašanje in kritično presodimo ugotovitve.

% =================================================
\section{Podatkovna analiza}\label{sec:podatki}
% =================================================

V tem poglavju s podatki za ZDA v obdobju 2010--2025 analiziramo vlaganja v R\&R in UI, investicije, nastajanje nove tehnologije in rast produktivnosti.

\subsection{Viri podatkov in izbor spremenljivk}

Vlaganja v R\&R merimo s skupnimi izdatki ZDA za R\&R v tekočih in stalnih cenah iz serije \emph{National Patterns of R\&D Resources}, ki jo NCSES objavlja v okviru poročila \emph{Science and Engineering Indicators} \parencite[tab.~DISC-1]{nsb2025}, ter z njihovim deležem v BDP. Za delež in za celotno stopnjo investicij uporabimo podatke BEA o nominalnem BDP in deležu bruto zasebnih domačih investicij v BDP \parencite{fredgdp}. Javna vlaganja v UI spremljamo s podatki programa NITRD \parencites{nitrd2021}{nitrd2025}, zasebna vlaganja in razširjenost uporabe UI s poročilom AI Index \parencite{aiindex2025}, nastajanje nove tehnologije pa s patentnimi prijavami, ki jih podatkovni niz USPTO AIPD prepozna kot povezane z UI \parencite{pairolero2025}. Produktivnost merimo s podatki U.S. Bureau of Labor Statistics (BLS) za zasebni nekmetijski poslovni sektor: z BDP na uro dela, skupno faktorsko produktivnostjo (TFP), rastjo realnega outputa, vložka kapitala in sestave delovne sile \parencites{fredtfp}{bls2026tfp}. Podatki o UI so redki in kratki, zato je empirični del predvsem analiza R\&R, investicij in produktivnosti, UI pa nastopa kot možen vzrok, ki ga preverjamo.

\begin{table}[H]
\centering
\caption{Kazalniki, uporabljeni v empirični analizi}\label{tab:kazalniki}
\small
\begin{tabularx}{\textwidth}{>{\raggedright\arraybackslash}p{3.0cm} >{\raggedright\arraybackslash}X >{\centering\arraybackslash}p{2.4cm} >{\centering\arraybackslash}p{2.3cm}}
\toprule
\textbf{Koncept} & \textbf{Kazalnik} & \textbf{Obdobje} & \textbf{Vir} \\
\midrule
Vlaganja v R\&R & Izdatki za R\&R v tekočih in stalnih cenah; delež v BDP & 2010--2023 & NSB/NCSES, BEA \\
Investicije & Delež bruto zasebnih domačih investicij v BDP & 2010--2025 & BEA \\
Vlaganja v UI & Federalni ne-obrambni izdatki za UI R\&R; zasebne naložbe v UI & FY2020--FY2025; 2024 & NITRD; AI Index \\
Nova tehnologija & UI-patentne prijave & 2018--2023 & USPTO \\
Produktivnost & BDP na uro dela, TFP, realni output, vložek kapitala, sestava delovne sile & 2010--2025 & BLS \\
\bottomrule
\end{tabularx}
\end{table}

\subsection{Metoda analize}

Empirični del temelji na opisni analizi časovnih vrst. Primerjamo podobdobja 2010--2014, 2015--2019 in 2020--2025, slednje pa razdelimo še na 2020--2022 in 2023--2025. Rast v podobdobju izračunamo kot povprečno letno (geometrijsko) stopnjo rasti med zadnjim letom prejšnjega in zadnjim letom tekočega podobdobja; prvo podobdobje se tako začne z izhodiščnim letom 2009. Ravni, kot sta delež R\&R ali investicij v BDP, prikazujemo kot povprečja let v podobdobju. Manjkajočih vrednosti ne interpoliramo, vsako serijo pa vzamemo iz ene same izdaje vira.

\subsection{Izdatki za raziskave in razvoj ter investicije}

Za izdatke za R\&R uporabimo serijo NCSES za obdobje 2010--2023 iz izdaje 2022--23 \parencite[tab.~DISC-1]{nsb2025}.

\begin{figure}[H]
\centering
\begin{minipage}{0.49\textwidth}
\centering
\begin{tikzpicture}
\begin{axis}[
    width=\linewidth,
    height=5.3cm,
    xlabel={Leto},
    ylabel={mrd USD},
    grid=major,
    xtick={2010,2013,2016,2019,2022},
    xticklabel={\pgfmathprintnumber[fixed,precision=0,1000 sep={}]{\tick}},
    tick label style={font=\footnotesize},
    label style={font=\footnotesize},
    legend pos=north west,
    legend style={font=\footnotesize}
]
\addplot+[mark=*,mark size=1.2pt] table[x=leto,y expr=\thisrow{skupaj_mio_usd}/1000,col sep=comma]{podatki_rnd_usa.csv};
\addlegendentry{tekoče cene}
\addplot+[mark=square*,mark size=1.2pt] table[x=leto,y expr=\thisrow{realno_mio_usd_2017}/1000,col sep=comma]{podatki_rnd_usa.csv};
\addlegendentry{stalne cene 2017}
\end{axis}
\end{tikzpicture}
\end{minipage}\hfill
\begin{minipage}{0.49\textwidth}
\centering
\begin{tikzpicture}
\begin{axis}[
    width=\linewidth,
    height=5.3cm,
    xlabel={Leto},
    ylabel={Delež v BDP (\%)},
    grid=major,
    xtick={2010,2013,2016,2019,2022},
    xticklabel={\pgfmathprintnumber[fixed,precision=0,1000 sep={}]{\tick}},
    yticklabel style={/pgf/number format/fixed,/pgf/number format/precision=1},
    tick label style={font=\footnotesize},
    label style={font=\footnotesize}
]
\addplot+[mark=*,mark size=1.2pt] table[x=leto,y=delez_bdp,col sep=comma]{podatki_rnd_usa.csv};
\end{axis}
\end{tikzpicture}
\end{minipage}
\caption{Izdatki za R\&R v ZDA, 2010--2023: levo v tekočih in stalnih cenah, desno kot delež BDP}
\caption*{\footnotesize Vir: \cite[tab.~DISC-1]{nsb2025}; \cite{fredgdp}; lastni izračun.}
\label{fig:rnd}
\end{figure}

Nominalni izdatki za R\&R so se povečali s 406,6 milijarde USD leta 2010 na 939,6 milijarde USD leta 2023, v stalnih cenah iz leta 2017 pa s 453,6 na 768,4 milijarde USD, kar je 4,1~\% na leto \parencite[tab.~DISC-1]{nsb2025}. Delež izdatkov za R\&R v BDP je bil do leta 2015 skoraj nespremenjen pri okoli 2,7~\%, kar ustreza vrednosti iz predavanj \parencite[pros.~131]{grobovsek2026}, nato pa je naraščal in leta 2023 dosegel 3,38~\% (lastni izračun). V povprečju podobdobij je znašal 2,70~\%, 2,86~\% in 3,37~\% (2020--2023). Večino ameriškega R\&R izvaja in financira poslovni sektor: leta 2023 je izvajal 78~\% in financiral 75~\% vseh izdatkov \parencite[tab.~DISC-1]{nsb2025}.

Delež bruto zasebnih domačih investicij v BDP je bil po recesiji nizek (13,3~\% leta 2009), nato pa je okreval in se ustalil pri okoli 18~\%. V povprečju podobdobij je znašal 16,0~\%, 17,8~\% in 18,1~\% \parencite{fredgdp}. Delež R\&R se je po letu 2015 torej povečal, celotna stopnja investicij pa skoraj ne.

\subsection{Federalna in zasebna vlaganja v UI}

NITRD objavlja federalna ne-obrambna vlaganja v UI R\&R in razlikuje med postavko \emph{Core UI}, ki zajema raziskave, pri katerih je UI neposredni predmet, in širšim seštevkom z dejavnostmi drugih področij s pomembno komponento UI \parencite{nitrd2025}.

\begin{table}[H]
\centering
\caption{Federalni ne-obrambni izdatki za UI R\&R}
\small
\begin{tabular}{lrr}
\toprule
\textbf{Fiskalno leto} & \textbf{Core UI (mio USD)} & \textbf{Skupaj (mio USD)} \\
\midrule
FY2020 dejansko & 797,5 & 1.426,0 \\
FY2021 odobreno & 954,9 & 1.534,4 \\
FY2023 dejansko & 1.914,2 & 3.121,9 \\
FY2024 odobreno & 1.789,0 & 2.977,5 \\
FY2025 zahteva & 1.954,5 & 3.316,1 \\
\bottomrule
\end{tabular}
\caption*{\footnotesize Vir: \cite{nitrd2021}; \cite{nitrd2025}; lastna priprava. Za FY2022 v uporabljenih virih ni primerljive vrednosti.}
\label{tab:UI-rd}
\end{table}

Dejanska poraba se je med FY2020 in FY2023 povečala s 1,43 na 3,12 milijarde USD, kar pa je le približno 0,3~\% vseh izdatkov za R\&R v ZDA leta 2023. Bistveno večja so zasebna vlaganja: zasebne naložbe v UI so v ZDA leta 2024 dosegle 109,1 milijarde USD, delež organizacij, ki poročajo o uporabi UI, pa se je med letoma 2023 in 2024 povečal s 55~\% na 78~\% \parencite{aiindex2025}. Zasebne naložbe niso enake izdatkom za R\&R, zato jih z drugimi serijami ne seštevamo.

\subsection{Nova tehnologija: UI-patenti}

Podatkovni niz USPTO AIPD za ameriške patentne dokumente z metodami strojnega učenja oceni verjetnost, da vsebujejo tehnologijo UI. Letno število UI-patentnih prijav v ZDA se je povečalo s približno 76.000 leta 2018 na nekaj več kot 101.000 leta 2023 (za 33~\%), prijave s področja strojnega učenja pa z okoli 12.000 na skoraj 29.000 (za 142~\%) \parencite{pairolero2025}. Patenti niso popolna mera inovacij, saj vse ideje niso patentirane, niz pa zajema tudi tuje prijavitelje. Patent je hkrati mehanizem izključljivosti, ki inovatorju omogoča monopolno rento in s tem spodbuja vlaganja v nove ideje \parencite[pros.~124--126]{grobovsek2026}.

\subsection{Razvoj produktivnosti}

BLS definira TFP kot razmerje med indeksom realnega outputa in indeksom kombiniranih vložkov dela in kapitala, pri čemer v kapital vključuje tudi intelektualno lastnino, med njo R\&R in programsko opremo \parencite{bls2026tfp}. Vložek kapitala BLS meri s kapitalskimi storitvami, ne z zalogo kapitala. Letne stopnje rasti izračunamo iz indeksov BLS iz izdaje marca 2026 \parencite{fredtfp}.

\begin{figure}[H]
\centering
\begin{tikzpicture}
\begin{axis}[
    width=0.94\textwidth,
    height=5.3cm,
    xlabel={Leto},
    ylabel={Rast TFP (\%)},
    grid=major,
    xtick=data,
    xticklabel={\pgfmathprintnumber[fixed,precision=0,1000 sep={}]{\tick}},
    xticklabel style={rotate=45,anchor=east},
    scaled y ticks=false
]
\addplot+[mark=*] table[x=leto,y=tfp_rast,col sep=comma]{podatki_tfp_usa.csv};
\end{axis}
\end{tikzpicture}
\caption{Letna rast TFP v zasebnem nekmetijskem poslovnem sektorju ZDA, 2010--2025}
\caption*{\footnotesize Vir: \cite{fredtfp}; lastni izračun.}
\label{fig:tfp}
\end{figure}

Gibanje TFP je precej bolj volatilno kot gibanje izdatkov za R\&R. Visoka rast leta 2010 (2,6~\%) je predvsem odboj po recesiji leta 2009, v letih 2011--2019 pa so se letne stopnje gibale med $-0{,}3$ in 1,3~\%. Pandemija je povzročila velika nihanja: po padcu leta 2020 je leta 2021 sledil skok za 3,9~\% in leta 2022 padec za 1,1~\%, v letih 2023--2025 pa je TFP rasla po 1,6~\%, 1,5~\% in 0,8~\%. BLS za poslovni cikel 2019--2025 poroča o 1,0-odstotni povprečni letni rasti TFP, za cikel 2007--2019 pa o 0,6-odstotni \parencite{bls2026tfp}. BDP na uro dela je v treh podobdobjih rasel po 1,10~\%, 1,44~\% in 2,23~\% na leto (lastni izračun iz indeksa) \parencite{fredtfp}. Kazalnik sestave delovne sile je leta 2020 poskočil za 2,8~\%, kar kaže na izrazito spremembo sestave zaposlenih med pandemijo \parencite{fredtfp}.

% =================================================
\section{Teoretični okvir: Solowov model}\label{sec:teorija}
% =================================================
Za analizo uporabimo Solowov model, ki temelji na prispevkih Solowa (1956) in Swana (1956) in uvaja minimalno teoretično strukturo, ki jo je mogoče empirično preveriti \parencite[pros.~12]{grobovsek2026}. Je temelj celotne teorije gospodarske rasti \parencite[pros.~39]{grobovsek2026} ter osrednji model rasti v zapiskih in učbeniku \parencites[str.~9--18]{polanec_rast}[pogl.~5--6]{garin2021}. Uporabimo različico z rastjo prebivalstva in tehnologije, saj osnovni Solowov model dolgoročne rasti ne predvideva \parencite[pros.~51]{grobovsek2026}. Dopolnimo jo z računovodsko razčlenitvijo rasti iz predavanj \parencite[pros.~64--81]{grobovsek2026}.

\subsection{Predpostavke modela}\label{sec:predpostavke}
\begin{itemize}
    \item \textbf{Diskreten čas in reprezentativni subjekti:} čas teče v obdobjih $t=1,2,3,\ldots$, gospodarstvo pa sestavljata reprezentativno podjetje in reprezentativno gospodinjstvo \parencite[pros.~15]{grobovsek2026}.
    \item \textbf{Neoklasična proizvodna funkcija:} je naraščajoča in konkavna v vsakem dejavniku, ima konstantne donose obsega ter izpolnjuje Inadine pogoje; poenostavimo jo s Cobb--Douglasovo obliko z delo povečujočo produktivnostjo $A$ \parencites[pros.~16]{grobovsek2026}[str.~10--12]{polanec_rast}.
    \item \textbf{Popolna konkurenca:} donosnost kapitala in plača sta enaki mejnima produktoma, delež kapitala v dohodku pa je enak $\alpha$ \parencite[pros.~17,~78]{grobovsek2026}.
    \item \textbf{Eksogena stopnja varčevanja:} gospodinjstvo ničesar ne optimizira, temveč varčuje eksogen delež $s\in(0,1)$ dohodka \parencites[pros.~18]{grobovsek2026}[str.~5]{garin2021}.
    \item \textbf{Zaprto gospodarstvo in neelastična ponudba dela:} investicije so enake varčevanju, ponudba dela je enaka povpraševanju \parencites[pros.~19,~21]{grobovsek2026}[str.~9]{polanec_rast}.
    \item \textbf{Eksogena rast prebivalstva in tehnologije:} delo raste po stopnji $n$, produktivnost po stopnji $g$, kapital se amortizira po stopnji $\delta$ \parencite[pros.~19,~52]{grobovsek2026}.
\end{itemize}

\subsection{Osnovne enačbe}\label{sec:enacbe}
Proizvodnja je
\begin{equation}\label{eq:prod}
Y_t=K_t^{\alpha}(A_tL_t)^{1-\alpha},
\end{equation}
kjer $A$ interpretiramo kot produktivnost, ki povečuje učinkovitost dela \parencite[pros.~16]{grobovsek2026}. Gospodinjstvo varčuje $S_t=sY_t$, investicije so enake prihrankom, $I_t=S_t$, kapital pa se kopiči po enačbi \parencite[pros.~19--21]{grobovsek2026}
\begin{equation}\label{eq:akum}
K_{t+1}=(1-\delta)K_t+sY_t .
\end{equation}
Prebivalstvo in tehnologija rasteta po konstantnih stopnjah, $L_t=L_1(1+n)^{t-1}$ in $A_t=A_1(1+g)^{t-1}$ \parencite[pros.~52]{grobovsek2026}. Sistem zapišemo s kapitalom na enoto učinkovitega dela $\tilde k_t\equiv K_t/(A_tL_t)$ \parencite[pros.~53--54]{grobovsek2026}:
\begin{equation}\label{eq:ktilde}
\tilde k_{t+1}=\frac{(1-\delta)\tilde k_t+s\,\tilde k_t^{\alpha}}{(1+g)(1+n)} .
\end{equation}
Ko je $\tilde k_t$ manjši od ustaljene vrednosti, $K$ in $Y$ rasteta hitreje kot $AL$; ko je večji, rasteta počasneje \parencite[pros.~56--57]{grobovsek2026}. Razlog je padajoči mejni produkt kapitala: vsaka dodatna enota kapitala ima vse manjši vpliv na proizvodnjo in investicije, dokler se investicije ne izenačijo z amortizacijo \parencite[pros.~31]{grobovsek2026}.

\subsection{Uravnotežena rast}\label{sec:bgp}
Na uravnoteženi poti rasti je $\tilde k_{t+1}=\tilde k_t=\tilde k^{\star}$. Iz enačbe~\eqref{eq:ktilde} sledi \parencite[pros.~55,~58]{grobovsek2026}
\begin{equation}\label{eq:bgp}
\tilde k^{\star}=\left(\frac{s}{g+n+gn+\delta}\right)^{\frac{1}{1-\alpha}},\qquad
y^{\star}_t=\frac{Y^{\star}_t}{L_t}=A_t\left(\frac{s}{g+n+gn+\delta}\right)^{\frac{\alpha}{1-\alpha}} .
\end{equation}
BDP na delavca torej dolgoročno raste s stopnjo tehnološkega napredka $g$; edini vir dolgoročne rasti dohodka na delavca je tehnologija \parencites[pros.~59]{grobovsek2026}[str.~13]{polanec_rast}. Stopnja varčevanja $s$, amortizacija $\delta$ in delež $\alpha$ določajo le dolgoročno \emph{raven} BDP na delavca; konvergenca zaradi kopičenja kapitala je prehoden pojav \parencites[pros.~61]{grobovsek2026}[str.~126--127]{garin2021}. Na uravnoteženi poti je razmerje med kapitalom in BDP konstantno in neodvisno od $A$ \parencite[pros.~77]{grobovsek2026}:
\begin{equation}\label{eq:ky}
\frac{K}{Y}=\frac{s}{g+n+gn+\delta}.
\end{equation}

\textbf{Primer.} Vzemimo parametre iz numeričnega primera na predavanjih, $\delta=0{,}06$, $\alpha=0{,}33$ in $s=0{,}2$ \parencite[pros.~40]{grobovsek2026}, ter $g=2~\%$ in $n=1~\%$, ki po zapiskih približno ustrezata ZDA \parencite[str.~18]{polanec_rast}. Po enačbi~\eqref{eq:ky} je $K/Y\approx2{,}2$ (lastni izračun), kar se ujema s Kaldorjevim dejstvom, da je razmerje skoraj konstantno in znaša okoli 2,2 \parencite[str.~9]{polanec_rast}. Učbenik za ZDA v obdobju 1950--2011 navaja razmerje med 3 in 3,5 \parencite[str.~61--62]{garin2021}; razlika izhaja iz različnih opredelitev kapitala in outputa, pomembna pa je skupna ugotovitev, da je razmerje dolgoročno približno konstantno.

\subsection{Računovodska razčlenitev rasti}\label{sec:racunovodstvo}
BDP je proizveden s funkcijo $Y=K^{\alpha}(ALh)^{1-\alpha}$, kjer je $h$ povprečni človeški kapital delavca. Na delavca jo lahko zapišemo na dva načina \parencite[pros.~65,~77]{grobovsek2026}:
\begin{equation}\label{eq:razclenitev}
y=A^{1-\alpha}k^{\alpha}h^{1-\alpha}
\qquad\text{oziroma}\qquad
y=A\left(\frac{k}{y}\right)^{\frac{\alpha}{1-\alpha}}h ,
\end{equation}
kjer sta $y=Y/L$ in $k=K/L$. Predavanje priporoča drugi zapis, ker kapital na delavca ni eksogen, saj je odvisen od $A$, razmerje $k/y=K/Y$ pa od $A$ ni odvisno \parencite[pros.~77]{grobovsek2026}. Z logaritmiranjem in časovno razliko dobimo (lastna izpeljava)
\begin{equation}\label{eq:rast}
g_y=g_A+\frac{\alpha}{1-\alpha}\,g_{K/Y}+g_h .
\end{equation}
Pri $\alpha=1/3$, kar ustreza povprečnemu dohodkovnemu deležu kapitala \parencites[pros.~78]{grobovsek2026}[str.~9,~23]{polanec_rast}, je $\alpha/(1-\alpha)=0{,}5$. Produktivnost $A$ izmerimo kot ostanek. Ker je $A^{1-\alpha}$ skupna faktorska produktivnost, velja $g_A\approx g_{\mathrm{TFP}}/(1-\alpha)$, kar uporabimo kot kontrolo z uradno TFP \parencites[pros.~16]{grobovsek2026}[str.~142]{garin2021}. Ostanek vsebuje tudi napake merjenja in vpliv drugih dejavnikov \parencites[str.~23--24]{polanec_rast}[str.~144--145]{garin2021}.

Razčlenitev ima eno res preverljivo napoved modela: na uravnoteženi poti je $g_{K/Y}=0$. Če se razmerje $K/Y$ ne spreminja, višje investicije ne prispevajo k rasti, ves preostanek pa je produktivnost in človeški kapital. Ker pa je $A$ ostanek in $g$ eksogen, model vsak pospešek, ki ni posledica investicij, pripiše tehnologiji, ne pove pa, od kod izvira.

\subsection{Uporaba modela na UI ter R\&R}\label{sec:uporabamodela}
UI v modelu ne nastopa neposredno, zato jo vključimo kot spremembo parametrov, podobno kot predavanje v numeričnih primerih spreminja $A$, $s$ in $L$ \parencites[pros.~45--48]{grobovsek2026}[str.~117--122]{garin2021}. Pri tem upoštevamo, da R\&R po teoriji ne služi kopičenju kapitala, temveč ustvarjanju idej in s tem $A$ \parencite[pros.~124]{grobovsek2026}, čeprav ga BLS v statistiki šteje med kapital \parencite{bls2026tfp}. Ločimo štiri možne razlage, ki imajo v razčlenitvi~\eqref{eq:rast} različen \enquote{podpis}:
\begin{itemize}
    \item \textbf{(A) Višja stopnja investicij $s$} (npr.\ naložbe v podatkovne centre in opremo za UI): višji $K/Y$ in raven BDP na delavca, rast pa le začasno \parencites[pros.~35,~47]{grobovsek2026}[str.~117--120]{garin2021}. Empirično merilo je celotna stopnja investicij $I/Y$, ne delež R\&R.
    \item \textbf{(B) Enkraten dvig ravni produktivnosti $A$:} BDP poskoči, kapital se prilagaja postopno; ob podvojitvi $A$ se BDP takoj poveča za faktor $2^{1-\alpha}\approx1{,}59$ \parencites[pros.~46]{grobovsek2026}[str.~120--122]{garin2021}. Če se nova tehnologija širi postopoma, se dvig ravni več let kaže kot višja rast $A$.
    \item \textbf{(C) Trajno višja stopnja rasti $g$:} BDP na delavca dolgoročno raste hitreje za $\Delta g$, $K/Y$ se po enačbi~\eqref{eq:ky} rahlo zniža. To je edina razlaga, ki trajno zviša rast.
    \item \textbf{(D) Višja intenzivnost R\&R v pol-endogenem modelu:} \textcite{jones1995} pokaže, da v modelu, v katerem rast nastaja z R\&R, dolgoročna stopnja rasti vseeno ni odvisna od politike, temveč le od parametrov, kot je rast prebivalstva; enako navaja predavanje \parencite[pros.~136]{grobovsek2026}. Iz tega sledi, da trajno višja intenzivnost R\&R prinese prehodno višjo rast $A$ in trajno višjo raven, ne pa trajno višje rasti. V razčlenitvi je podpis enak kot pri (B), vzrok pa ni UI, temveč višji delež R\&R, ki se je v ZDA povečal po letu 2015.
\end{itemize}

Učinke (A)--(C) smo številčno simulirali z enačbama~\eqref{eq:akum} in~\eqref{eq:prod} pri parametrih iz primera v razdelku~\ref{sec:bgp}. Gospodarstvo je do obdobja 0 na uravnoteženi poti, nato pa: (A) se $s$ poveča za eno odstotno točko; (B) se $A$ enkrat poveča za 3~\%; (C) se $g$ trajno poveča z 2 na 2,5~\%.

\begin{figure}[H]
\centering
\begin{tikzpicture}
\begin{axis}[
    width=0.88\textwidth,
    height=5.6cm,
    xlabel={Obdobje},
    ylabel={Rast BDP na delavca (\%)},
    grid=major,
    xmin=-2, xmax=30,
    ymin=1.9, ymax=2.6,
    legend pos=north east,
    legend style={font=\footnotesize},
    tick label style={font=\footnotesize},
    label style={font=\footnotesize},
    restrict y to domain=1.9:2.6
]
\addplot+[mark=none,thick,black,dashed] table[x=t,y=osnova,col sep=comma]{simulacija.csv};
\addlegendentry{osnova}
\addplot+[mark=none,thick,blue] table[x=t,y=A,col sep=comma]{simulacija.csv};
\addlegendentry{(A) $s$ +1 o.\,t.}
\addplot+[mark=none,thick,red] table[x=t,y=B,col sep=comma]{simulacija.csv};
\addlegendentry{(B) $A$ +3~\%}
\addplot+[mark=none,thick,teal] table[x=t,y=C,col sep=comma]{simulacija.csv};
\addlegendentry{(C) $g$ +0,5 o.\,t.}
\end{axis}
\end{tikzpicture}
\caption{Simulirana rast BDP na delavca v Solowovem modelu ob treh možnih učinkih}
\caption*{\footnotesize Opomba: v scenariju (B) rast v obdobju 0 znaša 4,04~\% in je zaradi preglednosti izven prikaza. Vir: lastni izračun; parametri iz \cite[pros.~40]{grobovsek2026} in \cite[str.~18]{polanec_rast}.}
\label{fig:simulacija}
\end{figure}

Pri (A) se rast BDP na delavca v prvem letu poveča za okoli 0,15 odstotne točke na odstotno točko višjega $s$, nato pa se postopoma vrne na 2~\%, $K/Y$ pa naraste z 2,22 proti 2,31. Pri (B) rast v obdobju dviga znaša 4,0~\%, nato pa je le malo nad osnovo. Pri (C) rast takoj naraste za $(1-\alpha)\cdot0{,}5\approx0{,}33$ odstotne točke in se nato približuje 2,5~\%, $K/Y$ pa pade proti 2,10 (lastni izračun). Tabela~\ref{tab:napovedi} povzema napovedi.

\begin{table}[H]
\centering
\caption{Napovedi za možne razlage v razčlenitvi~\eqref{eq:rast}}\label{tab:napovedi}
\small
\renewcommand{\arraystretch}{1.1}
\begin{tabularx}{\textwidth}{>{\raggedright\arraybackslash}p{4.0cm} >{\raggedright\arraybackslash}X >{\raggedright\arraybackslash}X >{\raggedright\arraybackslash}X}
\toprule
Razlaga & Prispevek $K/Y$ & Rast $A$ & Dolgoročna rast \\
\midrule
(A) višji $s$ & Začasno pozitiven & Nespremenjena & Nespremenjena \\
(B) enkraten dvig $A$ & Najprej negativen, nato pozitiven & Enkratno ali nekaj let višja & Nespremenjena \\
(C) trajno višji $g$ & Rahlo negativen & Trajno višja & Višja za $\Delta g$ \\
(D) višja intenzivnost R\&R & Kot pri (B) & Nekaj let višja, z zamikom za R\&R & Nespremenjena \\
\bottomrule
\end{tabularx}
\caption*{\footnotesize Vir: lastna priprava na podlagi \cite[pros.~45--48,~59--61,~136]{grobovsek2026}, \cite{jones1995} in slike~\ref{fig:simulacija}.}
\end{table}

% =================================================
\section{Empirična aplikacija Solowovega modela}\label{sec:uporaba}
% =================================================
Za vsako podobdobje rast BDP na uro dela razčlenimo po enačbi~\eqref{eq:rast} z $\alpha=1/3$, pri čemer je $L$ število opravljenih ur, $h$ pa kazalnik sestave delovne sile BLS. Ker BLS meri kapital s kapitalskimi storitvami, je $K/Y$ v naši razčlenitvi razmerje indeksov kapitalskih storitev in outputa, ne razmerje zaloge kapitala in BDP iz modela. Rast $K/Y$ izračunamo iz zaokroženih letnih sprememb. Model je dolgoročni model rasti, zato ne zajame recesije in pandemije; gospodarstvo navadno niha okoli svojega ustaljenega stanja \parencite[pros.~44,~49]{grobovsek2026}.

\begin{table}[H]
\centering
\caption{Razčlenitev rasti BDP na uro dela po enačbi~\eqref{eq:rast} (povprečna letna rast v \%)}\label{tab:obdobja}
\footnotesize
\renewcommand{\arraystretch}{1.15}
\begin{tabularx}{\textwidth}{>{\raggedright\arraybackslash}p{4.7cm} *{5}{>{\centering\arraybackslash}X}}
\toprule
 & 2010--2014 & 2015--2019 & 2020--2025 & \textit{2020--2022} & \textit{2023--2025} \\
\midrule
BDP na uro dela $g_y$ & 1,10 & 1,44 & 2,23 & 1,96 & 2,50 \\
\quad produktivnost $g_A$ & 1,22 & 1,10 & 1,71 & 1,18 & 2,25 \\
\quad razmerje $K/Y$: $0{,}5\,g_{K/Y}$ & $-0{,}48$ & $-0{,}06$ & $-0{,}05$ & 0,06 & $-0{,}15$ \\
\quad človeški kapital $g_h$ & 0,36 & 0,40 & 0,56 & 0,72 & 0,40 \\
\midrule
Kontrola: $g_{\mathrm{TFP}}/(1-\alpha)$ po BLS & 1,22 & 1,02 & 1,55 & 1,11 & 1,99 \\
Investicije $I/Y$ (povprečje) & 16,0 & 17,8 & 18,1 & 18,0 & 18,1 \\
Delež R\&R v BDP (povprečje) & 2,70 & 2,86 & 3,37\textsuperscript{a} & 3,36 & 3,38\textsuperscript{b} \\
\bottomrule
\end{tabularx}
\caption*{\footnotesize Opombe: $\alpha=1/3$. Rast je izračunana od zadnjega leta prejšnjega obdobja (npr. 2009--2014, 2022--2025). \textsuperscript{a}2020--2023. \textsuperscript{b}2023. Vir: \cite{fredtfp}; \cite{fredgdp}; slika~\ref{fig:rnd}; lastni izračun.}
\end{table}

Izmerjena produktivnost se z uradno TFP ujema v prvem podobdobju natančno, v drugih pa razlika ne presega 0,26 odstotne točke, ker BLS uporablja dejanske deleže stroškov kapitala. Edina preverljiva napoved modela se potrdi: po prvem podobdobju se je razmerje $K/Y$ ustalilo.

\subsection{Prvo obdobje: od 2010 do 2014}
\textbf{Uporaba modela:} BDP na uro je rasel po 1,10~\%, od tega 1,22 odstotne točke zaradi produktivnosti, 0,36 zaradi človeškega kapitala in $-0{,}48$ zaradi upadanja $K/Y$ za 0,96~\% na leto. Stopnja investicij je bila nizka (16,0~\%), delež R\&R pa stabilen.

\textbf{Primerjava z modelom:} upadanje $K/Y$ je skladno s prehodno dinamiko modela. Med recesijo je BDP padel bolj kot kapital, investicije pa so bile nizke, zato se je razmerje $K/Y$ vračalo proti uravnoteženi vrednosti \parencite[pros.~30,~56--57]{grobovsek2026}.

\textbf{Kritična presoja:} prispevek produktivnosti vključuje odboj leta 2010, ko je izmerjena $A$ zrasla za okoli 4~\%; brez tega leta bi bila rast $A$ v letih 2010--2014 le okoli 0,54~\% na leto (lastni izračun). Gre za ciklično izkoriščenost zmogljivosti, ki jo dolgoročni model ne pojasni.

\subsection{Drugo obdobje: od 2015 do 2019}
\textbf{Uporaba modela:} stopnja investicij je narasla na 17,8~\% (za 1,9 odstotne točke). Po simulaciji (A) bi to v prvem letu povečalo rast za okoli $1{,}9\cdot0{,}15\approx0{,}28$ odstotne točke (lastni izračun). BDP na uro je pospešil na 1,44~\% (za 0,34 odstotne točke), prispevek $K/Y$ pa se je izboljšal z $-0{,}48$ na $-0{,}06$, medtem ko se rast $A$ ni povečala (1,10~\%).

\textbf{Primerjava z modelom:} pospešek je po smeri in velikosti skladen z učinkom (A) oziroma s koncem prilagajanja po recesiji. Rast produktivnosti ni pospešila, čeprav je delež R\&R začel naraščati, zato znakov učinkov (C) ali (D) še ni.

\textbf{Kritična presoja:} nespremenjena rast produktivnosti kljub naraščajočim vlaganjem v R\&R je skladna z ugotovitvijo, da je rast produktivnosti v ZDA kljub hitremu povečevanju števila raziskovalcev ostala približno nespremenjena (Bloom in sod., 2020, povzeto po \cite[pros.~136--137]{grobovsek2026}), ter z dolgoročno upočasnitvijo rasti produktivnosti v ZDA, ki jo opisuje učbenik \parencite[str.~129]{garin2021}.

\subsection{Tretje obdobje: od 2020 do 2025}
\textbf{Uporaba modela:} stopnja investicij je ostala pri okoli 18~\% (le 0,2 odstotne točke več kot v drugem obdobju), zato učinek (A) po simulaciji prinese le okoli 0,03 odstotne točke dodatne rasti, $K/Y$ pa se ni povečal. BDP na uro je pospešil na 2,23~\%, rast $A$ na 1,71~\%. Po modelu je to mogoče le s hitrejšim tehnološkim napredkom, torej z razlagami (B), (C) ali (D).

\textbf{Robustnost:} pospešek ni enakomeren. V letih pandemije 2020--2022 je bila rast $A$ 1,18~\%, torej skoraj enaka kot v drugem obdobju, rast BDP na uro pa je napihnil prispevek sestave delovne sile (0,72), ki je leta 2020 poskočil za 2,8~\%. Brez leta 2020 bi bila rast BDP na uro v letih 2021--2025 le 1,58~\% s prispevkom človeškega kapitala 0,12 (lastni izračun). Pospešek produktivnosti je torej skoncentriran v letih 2023--2025, ko je rast $A$ znašala 2,25~\%, kar je za 1,15 odstotne točke več kot v drugem obdobju.

\textbf{Ali lahko pospešek pripišemo UI?} Tri ugotovitve kažejo, da zelo verjetno ne v celoti:
\begin{itemize}
    \item \textbf{Velikost:} \textcite[29]{acemoglu2024} ocenjuje, da bo UI letno rast TFP zvišala za okoli 0,06 odstotne točke, kar v enotah $A$ ustreza okoli 0,1 odstotne točke. To je manj kot desetina opaženega pospeška.
    \item \textbf{Časovni potek:} splošnonamenske tehnologije zahtevajo neopredmetene naložbe, ki jih statistika ne zajame, zato je izmerjena rast produktivnosti v začetku uvajanja podcenjena in šele pozneje precenjena \parencite[3]{brynjolfsson2018}. Generativna UI bi zato v letih 2023--2025 izmerjeno produktivnost prej znižala kot zvišala; takojšen skok po izidu orodja ChatGPT s tem ni skladen.
    \item \textbf{Alternativna razlaga:} delež R\&R v BDP se je povečal z 2,70~\% na 3,38~\%, predvsem po letu 2016. Po razlagi (D) bi to z zamikom prineslo prehodno višjo rast $A$ brez kakršnekoli vloge UI \parencite{jones1995}.
\end{itemize}

\textbf{Kritična presoja:} pospešek v letih 2023--2025 je skladen z (B), (C) in (D), s podatki pa jih ne moremo ločiti: tri leta in znižanje rasti TFP leta 2025 ne zadoščajo za presojo, ali je povečanje trajno. Ker je $A$ ostanek, vsebuje tudi napake merjenja in morebitne posledice prerazporeditve dela po pandemiji \parencite[str.~23--24]{polanec_rast}.

% =================================================
\section{Kritična presoja in zaključek}\label{sec:zakljucek}
% =================================================

\subsection{Odgovor na raziskovalno vprašanje}
Raziskovalno vprašanje je bilo: ali je pospešek rasti produktivnosti v ZDA v obdobju 2010--2025 po Solowovem modelu posledica višjih investicij ali hitrejšega tehnološkega napredka in ali ga lahko pripišemo umetni inteligenci?

(1) Vlaganja v R\&R in UI so se povečala: delež R\&R v BDP je narasel za 0,67 odstotne točke, federalna vlaganja v UI so se več kot podvojila, zasebne naložbe v UI pa so leta 2024 dosegle 109,1 milijarde USD. Celotna stopnja investicij se po letu 2015 ni povečala. (2) Preverljiva napoved modela se je potrdila: razmerje $K/Y$ se je po recesiji ustalilo, zato pospešek po letu 2019 ne izvira iz investicij, temveč iz produktivnosti, v letih pandemije pa tudi iz spremembe sestave delovne sile. (3) Pospešek produktivnosti je skoncentriran v letih 2023--2025. UI ga po velikosti in časovnem poteku verjetno ne pojasni v celoti: ocena njenega učinka je za red velikosti manjša, nove tehnologije pa izmerjeno produktivnost sprva prej znižajo. Enako dobro ga pojasni višja intenzivnost R\&R po letu 2015.

Na prvi del vprašanja je odgovor torej jasen: pospešek je posledica hitrejšega tehnološkega napredka, ne višjih investicij. Na drugi del pa podatki ne dajejo pritrdilnega odgovora: jasnega učinka UI na produktivnost v ZDA do leta 2025 še ni mogoče zaznati.

\subsection{Kaj model pojasni in česa ne}
Model dobro pojasni razliko med učinki na raven in na rast, prehodno upadanje $K/Y$ po recesiji in to, zakaj višje investicije ne morejo trajno zvišati rasti. Ne pojasni pa vira tehnološkega napredka, saj je $g$ eksogen \parencites[pros.~61]{grobovsek2026}[str.~18]{polanec_rast}; vsak pospešek, ki ni posledica investicij, model pripiše tehnologiji. To vprašanje obravnava teorija endogene rasti, v kateri ideje nastajajo z investicijami v R\&R, ki so odvisne od izključljivosti, njihov učinek pa krepijo tehnološki prelivi \parencite[pros.~122--128]{grobovsek2026}. Pol-endogeni model \parencite{jones1995} je za naše vprašanje posebej pomemben, ker pospešek pojasni brez UI. Solowov model predpostavlja tudi eksogeno stopnjo varčevanja \parencite[pros.~86]{grobovsek2026} in zaprto gospodarstvo, čeprav se znanje širi med državami \parencite[str.~31--32]{polanec_rast}.

\subsection{Omejitve in zaključek}
Analiza je opisna in temelji na kratkih podobdobjih. Razčlenitev predpostavlja stalen $\alpha=1/3$, produktivnost je ostanek z napakami merjenja in ciklično izkoriščenostjo zmogljivosti \parencite[str.~23--24]{polanec_rast}, $K/Y$ pa temelji na kapitalskih storitvah in zaokroženih podatkih. Delež R\&R v BDP velja za celotno gospodarstvo, podatki BLS pa za zasebni nekmetijski poslovni sektor. Podatki o UI so redki: zasebne naložbe imamo le za eno leto, patente od leta 2018. Korak od opisne analize proti identifikaciji bi bila primerjava panog, bolj in manj izpostavljenih UI, saj BLS objavlja produktivnost po panogah.

Solowov model pokaže, da UI dolgoročno rast lahko zviša le, če trajno pospeši tehnološki napredek, ne zgolj z višjimi investicijami. Podatki za ZDA kažejo hitrejšo rast produktivnosti v letih 2023--2025, vendar jo po velikosti in časovnem poteku težko pripišemo UI. Ali bo UI kot makro inovacija ovrgla oceno, da nam zmanjkuje idej \parencite[pros.~136--139]{grobovsek2026}, bo mogoče presoditi šele z daljšimi časovnimi vrstami.

\newpage
\printbibliography[title={Literatura in viri}, heading=bibintoc]

\end{document}